% --- CAPÍTULO 1: INTRODUÇÃO ---
\chapter{Introdução}
\label{cap:introducao}


Iniciei minha jornada nesta pesquisa buscando uma maneira de falar das contingências do autismo sem cair nos reducionismos característicos da prática científica moderna; 
    meu objetivo era bastante pessoal: 
        conseguir especular e descrever um mundo em que não seja inadequado ser \textit{assim}.

Ser \textit{assim} é habitar cada momento sem apagar as contingências;
    por algum tempo eu entendia que ser autista seria como ser "impermeável à cultura", que ela não nos penetrasse, passasse direto, mas o percurso intelectual que percorri desde aquele primeiro passo me trouxe a outra compreensão:
        ser \textit{assim} é ter um corpo mais aberto a diferenças, um corpo que literalmente olha para fora antes de olhar para dentro.

Uma IA treinada para diagnosticar autismo em bebês recém-nascidos rastreando o percurso de seus olhos ao serem apresentados a diferentes imagens tem demonstrado ser excepcionalmente bem sucedida;
    há um padrão bem marcante que determina o diagnóstico:
        diante de uma cena, os bebês autistas percorrem com o olhar todo o cenário antes de direcioná-loo aos atores, ao contrário dos bebês não diagnosticados, que rapidamente direcionam o olhar para os agentes em cena.

Melhor dizendo --- e, talvez, a transposição discursiva que farei aqui seja emblemática do que tenho a dizer --- os bebês autistas orientam seu olhar a uma gama mais vasta de agentes;
    entre os não diagnosticados, muito do que está agindo e afetando seus colegas é registrado como cenário, como contingência, não como agência:
        ocasionalmente, uma especulação emergiu da consciência de que, naquele mundo que almejo descrever, não há lugar para a visão objetiva que orienta a prática e o discurso científicos na modernidade.

E depois de muito resistir à definição de uma pergunta de pesquisa para o meu trabalho;
    e de tanto sofrer para falar a seu respeito -- seja para um colega ou professor, seja para um familiar ou amigo -- notei que havia efetivamente uma pergunta me orientando, com a qual eu sempre estive me debatendo, e que foi um ponto de controvérsia que se colocou a cada passo da pesquisa, a cada trabalho ou comunicação a ser realizado, demarcando ao longo do tempo toda uma superfície de contato entre mim e a interlocução que encontrei sempre tão inacessível:
        afinal, \textbf{como falar cientificamente?}

Sem conseguir escrever minhas ideias de um modo inteligível a interlocutor algum, sem saber como eu as poderia organizar numa dissertação de mestrado, eu também não sabia como deveria realizar minha pesquisa, com quais métodos ou teorias, conceitos ou categorias, articularia um discurso científico que desarticula o discurso científico da modernidade;
    parece uma inversão lógica orientar a pesquisa pela comunicação que se fará dela, e várias outras inversões devem aparecer ao longo desta leitura: 
        esta inversão é fundamental para compreender minha proposição por um olhar contra-objetivista para orientar o discurso -- e, portanto, a prática -- científico.

Contra-objetivismo é uma tentativa de construir uma associação com as palavras de Nego Bispo, cujo livro "A terra dá, a terra quer" tive o prazer de reler apenas uma semana antes de iniciar esta escrita, anos depois de minha primeira leitura;
    e neste reencontro descobri a mediação contingente à presente com-posição:
        penso-com Nego Bispo para propor como ficar com o problema da objetividade no discurso científico.

Objetivamente, o que a diferença de olhar entre os bebês examinados pela IA mencionanda anteriormente me sensibiliza a notar é a objetividade atuando como um artifício do corpo para habitar um ambiente;
    ela não articula o mundo como ele é, e nem atua do mesmo modo entre os bebês diagnosticados e os não diagnosticados:
        a objetividade é diferencial de um modo que jamais aparece ao olhar objetivista, ela é posicional e se dispõe em camadas.

O olhar que agencia o cenário é externo àquele que o objetifica, mas não superior ou mais amplo;
    ao menos não apenas por lhe ser externo:
        estou proponto que o ponto cego do olhar objetivista seja em muito mediado por sua negação da possibilidade de constituir ou comunicar ciência fora dos limites da objetividade, negação que alcança as agências dos cenários, dos territórios, da natureza e, até mesmo, dos próprios sujeitos -- inclusive daqueles que, a fim de enxergarem objetivamente, se fazem de objetos.

Objetivismo e colonialismo rimam na fonética e na semântica; 
    ambos manifestam a mesma terceira e cega rima:
        o monoteísmo.

Esta palavra, em dança com Nego Bispo, agencia uma encruzilhada oportuna nos olhares objetivistas;
    quando Bispo fala de monoteísmo, escuto-o compartilhar uma denúncia crítica ao higienismo coloialista:
        entre seus críticos objetivistas, a leitura é sempre associada ao sentido usual da palavra, remetendo à associação religiosa a um único ser divino.

Em vez de contrapor as leituras, irei sobrepô-las, como Velasquez em Las Meninas;
    e utilizarei a heterotopografia como método:
        Foucault propõe uma ciência das heterotopias, uma forma especial que, dentre outras características, o francês descreve sobrepor elementos que usualmente não fazem sentido postos no mesmo espaço.

A heterotopografia é contra-objetivista porque toma proposições objetivamente contraditórias para compor com elas um sentido comum;
    o olhar objetivista não comporta uma construção na qual duas leituras distintas e opostas de um mesmo fenômeno componham, sem negarem uma a outra, um discurso articulado e científico:
        para este olhar, a contradição se associa a uma limitação do discurso, para o contra-objetivismo, a uma limitação da linguagem científica.
        
Uma palavra, ou um signo, é uma seta;
    um significante que aponta para um significado:
        o objetivismo é um olhar bastante eficiente em concatenar setas, em organizá-las em fileiras ordenadas, em dividi-las, categorizá-las.

No contra-objetivismo, a palavra é duas setas;
    uma que chega, e outra que sai:
        o significado ativo apontando para o significante, e este para o significado passivo.

Proponho a nomeação como um processo de sobreposição semântico-semiótica no qual um significante constitui uma sinergia entre um nomeado -- significado passivo -- e um nome -- significado ativo --;
    a seta sai do nome e chega no significante, e então vai deste para o nomeado:
        assim, se construo um objeto com a intenção de construir uma cadeira, o objeto é o nomeado, "cadeira" é o nome, e o significante sou eu.

No entanto, nomear tem seus custos;
    a intenção paga um preço para informar o nomeado:
        o significado passivo o é apenas retoricamente, pois suas contingências e recalcitrâncias possuem agência no processo de informação também. elas delimitam as possibilidades de formatação do objeto segundo o seu nome, nomeando-o de volta, isto é, informando-o segundo aquilo que o nome informa por origem.

Nomear justapõe dois significados distintos e constitui uma sinergia informacional entre eles, pela qual um dá forma ao outro de modo recíproco e contínuo, mas não simétrico;
    o significado ativo, o nome, é anterior ao passivo, é seu ancestral, e dá forma ao nomeado tal qual o sobrenome de um bebê recém nascido lhe dá sua identidade:
        mas, assim como aquele chega com suas contingências mediando tanto a identidade quanto o nome que recebe de sua família, a madeira de que é feito o objeto, os pregos, as ferramentas e técnicas utilizadas, e quaisquer outras contingências do processo que lhe dá forma de cadeira também determinam sua própria forma, e contra-informam aquela da cadeira.

O olhar monoteísta visualiza o objeto e só enxerga uma cadeira;
    o politeísta enxerga naquele a cadeira, mas não o reduz a isto:
        uma cadeira pode ser um anteparo que bloqueie uma porta, ou uma escada que auxilie a alcançar alturas maiores, pode ser uma arma branca, uma herança de família, um presente do território, uma ferida contra este, pode ser, e é, muitas coisas, pois tudo que lhe for contingente lhe informa e é por ela informado.

Proponho que falar cientificamente envolve localizar na fala tanto o significante quanto ambos os significados, passivo e ativo, acolher articulações distintas que mobilizem um ou mais destes elementos, e não pretender jamais afirmar ou esgotar nada do que se diz;
    entrelaço assim a ideia de \textit{proposições bem articuladas} de Latour, e a de \textit{objetividade parcial e localizada} de Haraway, para compor uma gramática científica que reoriente não apenas o modo como se fala cientificamente, mas também o modo como se pratica ciência:
        para ficar com o problema da objetividade, proponho deslocar o olhar para fora e contranger a objetividade para dentro deste olhar, para o significante que articula as setas entre os significados.

Proponho uma metaforma que é um conjunto infinito de metaformas;
    chamo esta metaforma de Totem:
        o Totem é esta metaforma que articula significante, significado passivo e significado ativo.

Proponho figurar o significado passivo como profano, e o ativo como sagrado;
    o significante, por sua vez, figura o hierofante, o sacerdote que medeia o divino e o mundano:
        o Totem como uma metaforma da religiosidade.
    
Proponho figurar o significado passivo como início e o ativo como fim;
    o significante, como esperado, figura o meio:
        o Totem como uma metaforma da linearidade.
    
Proponho figurar o significado passivo \textit{\textbf{e}} o ativo como começo; 
    o significante, como meio:
        o Totem como uma metaforma da confluência e do tempo espiralar.

Proponho o Totem como uma metaforma hierofante mediadora de mundos;
    mundos tão controversos quanto o são o sagrado e o profano entre si, constituindo línguas comuns:
        o Totem como um artifício para ficar com o problema.

Proponho uma prática científica que articule a especulação com a proposição, e que abandone as hipóteses;
    especulação como a que fazem os morcegos ao gritarem para se orientar no mundo:
        primeiro orientar-se, especular um mundo específico em torno de um propósito, e depois, propor, articular os seus próprios ecos a algum território.

Proponho a prática de entrar no bosque da vida coletando frutos e estórias, para então nos sentarmos numa clareira no centro da floresta, estendermos uma longa toalha no chão, e compartilharmos os achados trazidos por cada ser num grande piquenique;
    fazer ciência como quem enche uma bolsa de ficções científicas para trocar com outros colecionadores de ficção científica:
        afinal de contas, \textit{objetivamente falando, todo saber científico não passa de um belo faz-de-conta}.
    
Não que isto seja um problema;
    pelo contrário:
        a prática contra-objetivista \textit{alça} a ciência ao patamar da poesia.


---

Navegava certa vez por uma rede social qualquer quando tropecei no relato sobre um professor que orientava seus alunos a lrem como se mirassem o mar e, ocasionalmente, avistassem baleias emergindo e saltando através de sua superfície ondulante;
    sem qualquer preocupação em enxergar todos os seus mistérios, mas atentos àqueles com os quais se deparassem:
        orientação que reconheço como uma das mais necessárias para a leitura de textos e mundos que venho realizando desde então.

Mais recentemente, li uma carta de Bruno Latour -- falecido pensador francês das técnicas e ciências com o qual me associo na tarefa de \textit{composição de um \textbf{mundo comum}} -- endereçada a uma indefinida \textit{senhorita} -- aparentemente uma aluna, quicá uma orientanda -- inquieta e confusa quanto ao curso do professor Latour;
    ele tenta apaziguá-la elaborando os assuntos, a abordagem e a proposta do curso, mas admite ser necessário "\textit{ir até o fim para compreender do que se trata}":
        um jeito certamente curioso de apaziguar o coração de uma estudante cheia de dúvidas no qual reconheci, entretanto, uma das orientações mais necessárias para me acompanhar na escrita que inicio.

Ainda mais recentemente, li um relato de Givânia Maria da Silva -- quilombola educadora do quilombo de Conceição das Crioulas -- sobre a história transmitida oralmente de seu território;
    ela conta das mulheres que chegaram ali no ano de 1802 e compraram aquelas terras com o fruto de seu trabalho fiando algodão para vender na cidade, e que um homem chegou depois trazendo uma imagem de Nossa Senhora da Conceição na intenção de convertê-las ao Catolicismo, mas que "elas trouxeram a santa para o lado delas, houve um pacto feminino com a santa, e o homem perde a importância naquela história, ninguém sabe mais o que aconteceu":
        reconheço aqui uma sabedoria necessária, feminista e poderosa para aprender como ficar com o problema.

O problema com o qual escolhi ficar na pesquisa que relato e realizo neste texto é o \textbf{objetivismo};
    este é o nome que proponho para articular o interesse pessoal que me remeteu a este trabalho com o interesse coletivo necessário à realização de uma pesquisa científica:
        ao menos segundo a ética e gramática científicas que almejo propor.

Minha jornada teve início no dia dois de abril de 2024;
    até então, vinha tentando ingressar num programa de mestrado para fazer uma leitura do romance "Viva o povo brasileiro", de João Ubaldo Ribeiro, que discutisse o papel das construções simbólicas e históricas para a configuração social do Brasil contemporâneo:
        contudo, naquela data -- Dia Mundial de Conscientização do Autismo -- a publicação de uma página de divulgação científica criticando a expressão pública de pessoas autistas acerca de nossas experiências e reduzindo o diagnóstico a a uma confluência com a ordem neoliberal despertou em mim a necessidade de encontrar uma maneira de falar sobre as contingências do autismo sem cair nos reducionismos característicos da prática científica moderna.

Meu objetivo era -- e continua sendo -- conseguir especular e descrever um mundo em que não seja inadequado ser \textit{assim};
    um objetivo bastante pessoal para uma pesquisa científica, exatamente o que Latour propõe em um dos critérios de seu "Princípio de Falseabilidade de Stengers-Despret" para guiar uma ciência propositiva e capaz de fomentar a composição de um mundo comum:
        quanto mais próximo e interessado estiver o investigador científico de seu problema de pesquisa, mais e melhor ele lhe oferecerá escuta.

Mas eu preciso de uma baleia para o público não autista se espero que ele acompanhe a leitura até o final para compreender do que estou tratando;
    para falar, eu preciso ser ouvido, e para ser ouvido, eu preciso trazer o público para o meu lado, e para isso proponho um pacto contro o objetivismo:
        um \textbf{pacto contra-objetivista}.

Neste pacto, caberá a você que me lê prosseguir com atenção não só às baleias, mas também às controvérsias, contradições, paradoxos, ambiguidades, incompletudes e incoerências;
    quanto a mim, fico comprometido a conduzir-lhes por \textbf{uma jornada em busca de uma nova bússola para comunicações -- e práticas -- científicas bem articuladas}:
        associo assim ao sucesso de meu objetivo a proposição de uma gramática científica -- como diria Nego Bispo -- \textbf{politeísta}.

Se por um lado faço o melhor que posso nesta introdução para apresentar um ponto de partida inteligível, por outro carrego comigo o mesmo problema de Latour com seu curso;
    deixo desde já advertido que este é um texto espiralar, não linear:
    


% O ato de \textit{falar cientificamente} transcende a mera aplicação de um léxico especializado; ele constitui a adoção de um regime de visibilidade, validação e autoridade discursiva. Esta dissertação investiga os limites e as configurações desse discurso a partir do tensionamento provocado pelas formulações conceituais e cosmológicas de Antônio Bispo dos Santos, o Nego Bispo.

% \section{Problema de Pesquisa}
% Como se estruturam as fronteiras do discurso científico quando confrontadas por matrizes de conhecimento contracoloniais e baseadas na oralidade? A partir do debate contemporâneo em torno da recepção acadêmica de Nego Bispo, questiona-se de que maneira os critérios de cientificidade operam mecanismos de inclusão, tradução ou epistemicídio.

% \section{Objetivos}
% O objetivo geral desta pesquisa é analisar a engenharia do discurso de autoridade científica e suas clivagens diante de saberes orgânicos e afrodiaspóricos.